\section{The Conjecture for Three Agents}
\label{sec:n3}


\begin{theorem}
\label{thm:n3-main}
For every instance with three agents and negative dichotomous valuations, there
is an envy-free solution $(\alloc, \subsidy)$ with $\subsidy \in \set{0,1}^3$,
hence with total subsidy at most $2$, computable in polynomial time given
value-oracle access.
\end{theorem}


\begin{theorem}[\cite{TWYZ23}]
\label{thm:n3-twyz}
For any number $n$ of agents and any dichotomous cost functions there is a
polynomial-time algorithm producing a partial allocation
$(\bundle1,\dots,\bundle n)$ that is envy-free with $\subsidy = 0$ and leaves
at most $n-1$ chores unassigned.
\end{theorem}

At $n = 3$ this yields an envy-free partial allocation together with a residue
$R := \items \setminus (\bundle1 \cup \bundle2 \cup \bundle3)$ of at most two
chores. The case $R = \emptyset$ is trivial.
In the next sections we will prove the cases for $ \abs R = 1 $ and $\abs R = 2$

% --------------------------------------------------------------------------
\subsection{Excess}
\label{sec:n3-excess}

\begin{definition}[Excess]
\label{def:n3-excess}
Fix bundles $Y_1,\dots,Y_n$ and an assignment $\sigma \in \mathbb{S}_n$, agent
$a$ receiving $Y_{\sigma(a)}$. Put $\beta_a := \min_{j} \cost_a(Y_j)$ and
\[
  \varepsilon_a(\sigma) \;:=\; \cost_a(Y_{\sigma(a)}) - \beta_a \;\ge\; 0 .
\]
The \emph{total excess} of $\sigma$ is $\sum_{a \in \agents}
\varepsilon_a(\sigma)$.
\end{definition}

\begin{lemma}
\label{lem:n3-pathbound}
Let $\alloc$ be the allocation assigning $Y_{\sigma(a)}$ to $a$. Every simple
directed path in $\envygraph{\alloc}$ has weight at most $\sum_a
\varepsilon_a(\sigma)$; hence $\pathw{\alloc}(i) \le \sum_a
\varepsilon_a(\sigma)$ for every $i \in \agents$.
\end{lemma}

\begin{proof}
Each arc leaving $a$ has weight
$\edgew{\alloc}(a,b) = \cost_a(Y_{\sigma(a)}) - \cost_a(Y_{\sigma(b)}) \le
\cost_a(Y_{\sigma(a)}) - \beta_a = \varepsilon_a(\sigma)$, because $\beta_a$ is
a minimum over all of $Y_1,\dots,Y_n$. A simple path leaves each agent at most
once, and the empty path has weight $0$.
\end{proof}

\begin{lemma}
\label{lem:n3-mincost}
For every assignment $\sigma$,
$\sum_a \cost_a(Y_{\sigma(a)}) = \sum_a \beta_a + \sum_a
\varepsilon_a(\sigma)$. Since $\sum_a \beta_a$ does not depend on $\sigma$, an
assignment minimises total cost if and only if it minimises total excess.
\end{lemma}

\begin{proof}
Sum Definition~\ref{def:n3-excess} over $a \in \agents$.
\end{proof}

\begin{corollary}
\label{cor:n3-criterion}
Let $Y_1,\dots,Y_n$ be bundles admitting some assignment of total excess at
most $1$, and let $\pi$ minimise $\sum_a \cost_a(Y_{\pi(a)})$. Then the
allocation $\alloc$ assigning $Y_{\pi(a)}$ to $a$ admits an envy-free solution
with $\optsubsidy \in \set{0,1}^n$.
\end{corollary}

\begin{proof}
By Lemma~\ref{lem:n3-mincost}, $\pi$ has total excess at most $1$. In cost form
$\sum_a v_a(Y_{\pi(a)}) = -\sum_a \cost_a(Y_{\pi(a)})$, so minimising total
cost across reassignments of $Y_1,\dots,Y_n$ is exactly
condition~(ii) of Theorem~\ref{thm:hs-characterisation}; hence $\alloc$ is
envy-freeable. Its minimal subsidy is $\optsubsidy_i = \pathw{\alloc}(i)$ by
Theorem~\ref{thm:hs-minsubsidy}, which is at most $1$ by
Lemma~\ref{lem:n3-pathbound}, and integral and non-negative by
Observation~\ref{obs:integrality}.
\end{proof}

% --------------------------------------------------------------------------
\subsection{One unallocated chore}
\label{sec:n3-one}

The following holds for every number of agents; only the case $n = 3$ is needed
for Theorem~\ref{thm:n3-main}.

\begin{theorem}
\label{thm:n3-one}
Let $(\bundle1,\dots,\bundle n)$ be an envy-free partial allocation leaving
exactly one chore $e$ unassigned. Fix any $m \in \agents$, put
$Y_m := \bundle m \cup \set{e}$ and $Y_j := \bundle j$ for $j \ne m$, and let
$\pi$ minimise $\sum_a \cost_a(Y_{\pi(a)})$. Then the allocation assigning
$Y_{\pi(a)}$ to $a$ admits an envy-free solution with
$\optsubsidy \in \set{0,1}^n$.
\end{theorem}

\begin{proof}
By Corollary~\ref{cor:n3-criterion} it suffices to exhibit one assignment of
total excess at most $1$; the identity assignment is one.

Let $a \ne m$. For every $j$,
\[
  \cost_a(Y_a) \;=\; \cost_a(\bundle a) \;\le\; \cost_a(\bundle j)
  \;\le\; \cost_a(Y_j),
\]
the first inequality by envy-freeness of the partial allocation and the second
by monotonicity, since $\bundle j \subseteq Y_j$. Hence $\beta_a =
\cost_a(Y_a)$ and $\varepsilon_a(\mathrm{id}) = 0$.

For $a = m$ and every $j$,
\[
  \cost_m(Y_m) \;=\; \cost_m(\bundle m) + \cost_m(e \mid \bundle m)
  \;\le\; \cost_m(\bundle m) + 1
  \;\le\; \cost_m(\bundle j) + 1
  \;\le\; \cost_m(Y_j) + 1,
\]
using Definition~\ref{def:dichcost}, then envy-freeness, then monotonicity.
Taking the minimum over $j$ gives $\cost_m(Y_m) \le \beta_m + 1$, that is
$\varepsilon_m(\mathrm{id}) \le 1$.
\end{proof}

% --------------------------------------------------------------------------
\subsection{Two unallocated chores}
\label{sec:n3-two}

% This case uses more than the statement of Theorem~\ref{thm:n3-twyz}: it uses
% the state in which its algorithm halts. That algorithm maintains an envy-free
% partial allocation $(\bundle1,\dots,\bundle n)$ together with the
% \emph{equality graph}
% \begin{equation}
% \label{eq:n3-equalitygraph}
%   G = (\agents, E), \qquad
%   (i,j) \in E \iff i \ne j \ \text{ and } \ \cost_i(\bundle i) = \cost_i(\bundle j),
% \end{equation}
% and repeats three rules in the order below, halting only when none applies.
% \begin{enumerate}
%   \item[(R1)] If some unassigned chore $e$ and some agent $i$ satisfy
%         $\cost_i(e \mid \bundle i) = 0$, assign $e$ to $i$.
%   \item[(R2)] If some unassigned chore $e$ and some arc $(i,j) \in E$ lying on
%         a directed cycle of $G$ satisfy $\cost_i(e \mid \bundle j) = 0$, rotate
%         the bundles along that cycle and assign $e$ to $i$.
%   \item[(R3)] Otherwise select a strongly connected component $S$ of $G$ with
%         no arc of $E$ leaving $S$. If at least $\abs S$ chores are unassigned,
%         assign one to each agent of $S$; otherwise halt.
% \end{enumerate}

\begin{lemma}
\label{lem:n3-stopping}
Suppose the algorithm of Theorem~\ref{thm:n3-twyz}, run on three agents, halts
with partial allocation $(\bundle1,\bundle2,\bundle3)$ and residue $R$ of size
$2$. Then, with $G$ as in \eqref{eq:n3-equalitygraph}:
\begin{enumerate}
  \item[(a)] $G$ is strongly connected;
  \item[(b)] $\cost_i(e \mid \bundle i) = 1$ for every $i \in \agents$ and
        every $e \in R$;
  \item[(c)] $\cost_i(e \mid \bundle j) = 1$ for every arc $(i,j) \in E$ and
        every $e \in R$.
\end{enumerate}
\end{lemma}

\begin{proof}
Chores remain unassigned, so the algorithm halted at the final clause of (R3);
in particular neither (R1) nor (R2) applied at the halting state.

(b) If $\cost_i(e \mid \bundle i) = 0$ for some $i$ and some $e \in R$, rule
(R1) would apply. Marginals lie in $\set{0,1}$, hence the value is $1$.

(a) The component $S$ selected by (R3) satisfies $\abs R < \abs S$. As
$\abs R = 2$ this forces $\abs S = 3$, so $S = \agents$; and $S$ is strongly
connected by construction.

(c) Let $(i,j) \in E$. By (a) there is a directed path in $G$ from $j$ to $i$;
choosing one with no repeated vertex and appending the arc $(i,j)$ produces a
directed cycle of $G$ through $(i,j)$. If $\cost_i(e \mid \bundle j) = 0$ for
some $e \in R$, rule (R2) would apply to that arc and cycle.
\end{proof}

\begin{lemma}
\label{lem:n3-expensive}
In the situation of Lemma~\ref{lem:n3-stopping}, for every $a, j \in \agents$
and every $e \in R$,
\[
  \cost_a(\bundle j \cup \set{e}) \;\ge\; \cost_a(\bundle a) + 1 .
\]
\end{lemma}

\begin{proof}
Write $\cost_a(\bundle j \cup \set{e}) = \cost_a(\bundle j) +
\cost_a(e \mid \bundle j)$. Envy-freeness of the partial allocation gives
$\cost_a(\bundle j) \ge \cost_a(\bundle a)$, and both sides are integers by
Observation~\ref{obs:integrality}. If that inequality is strict the claim
follows, marginals being non-negative. If it is an equality then either
$j = a$, and $\cost_a(e \mid \bundle a) = 1$ by
Lemma~\ref{lem:n3-stopping}(b); or $j \ne a$, in which case $(a,j) \in E$ by
\eqref{eq:n3-equalitygraph} and $\cost_a(e \mid \bundle j) = 1$ by
Lemma~\ref{lem:n3-stopping}(c).
\end{proof}

\begin{theorem}
\label{thm:n3-two}
Let $(\bundle1,\bundle2,\bundle3)$ be the partial allocation at which the
algorithm of Theorem~\ref{thm:n3-twyz} halts with residue
$R = \set{e_1,e_2}$. Assign $e_1$ and $e_2$ to two \emph{distinct} bundles, and
assign the three resulting bundles to the three agents by a minimum-cost
matching. The resulting allocation admits an envy-free solution with
$\optsubsidy \in \set{0,1}^3$.
\end{theorem}

\begin{proof}
Write $Y_1,Y_2,Y_3$ for the bundles after placement: two are \emph{augmented},
carrying one chore of $R$ each, and one is \emph{clean}. Put
$L := \sum_a \cost_a(\bundle a)$ and, for an assignment $\sigma$,
\[
  g_a(\sigma) \;:=\; \cost_a(Y_{\sigma(a)}) - \cost_a(\bundle a) \;\ge\; 0,
\]
non-negative because $\cost_a(Y_{\sigma(a)}) \ge \cost_a(\bundle{\sigma(a)})
\ge \cost_a(\bundle a)$ by monotonicity and envy-freeness of the partial
allocation.

\emph{The minimum of $\sum_a \cost_a(Y_{\sigma(a)})$ is exactly $L + 2$.} For
the lower bound, fix $\sigma$: the two agents receiving augmented bundles have
$g_a(\sigma) \ge 1$ by Lemma~\ref{lem:n3-expensive}, and the third has
$g_a(\sigma) \ge 0$. For the upper bound, hand each bundle to its original
owner: the two agents whose own bundles were augmented pay
$\cost_a(\bundle a) + 1$ by Lemma~\ref{lem:n3-stopping}(b), and the third pays
$\cost_a(\bundle a)$.

\emph{Hence at a minimum-cost assignment $\pi$ the excesses are pinned:}
$\sum_a g_a(\pi) = 2$ with $g_a(\pi) \ge 1$ on each of the two
augmented-bundle holders, so both equal $1$ exactly and the agent $c$ holding
the clean bundle has $g_c(\pi) = 0$.

Let $a$ and $b$ be the augmented-bundle holders, and recall
$\edgew{}(x,y) = \cost_x(Y_{\pi(x)}) - \cost_x(Y_{\pi(y)})$.
\begin{itemize}
  \item $\edgew{}(a,b) = \bigl(\cost_a(\bundle a) + 1\bigr) -
        \cost_a(Y_{\pi(b)}) \le 0$: the first term because $g_a(\pi) = 1$, and
        $\cost_a(Y_{\pi(b)}) \ge \cost_a(\bundle a) + 1$ by
        Lemma~\ref{lem:n3-expensive}, the bundle $Y_{\pi(b)}$ being augmented.
        Symmetrically $\edgew{}(b,a) \le 0$.
  \item $\edgew{}(a,c) = \bigl(\cost_a(\bundle a) + 1\bigr) -
        \cost_a(Y_{\pi(c)}) \le 1$, since $Y_{\pi(c)}$ is clean and therefore
        costs $a$ at least $\cost_a(\bundle a)$. Symmetrically
        $\edgew{}(b,c) \le 1$.
  \item $\edgew{}(c,x) = \cost_c(\bundle c) - \cost_c(Y_{\pi(x)}) \le 0$ for
        both $x$, since $g_c(\pi) = 0$ gives the first term and every bundle
        costs $c$ at least $\cost_c(\bundle c)$.
\end{itemize}
Every arc therefore has weight at most $0$ except the two arcs entering $c$,
which have weight at most $1$. A simple path enters $c$ at most once, every arc
leaving $c$ has weight at most $0$, and every other arc available to it has
weight at most $0$; hence no simple path has weight above $1$. The assignment
is minimum-cost, so the allocation is envy-freeable by
Theorem~\ref{thm:hs-characterisation}, and its minimal subsidy is
$\optsubsidy_i = \pathw{}(i) \le 1$ by Theorem~\ref{thm:hs-minsubsidy},
integral and non-negative by Observation~\ref{obs:integrality}.
\end{proof}

\begin{remark}
\label{rem:n3-construction}
Theorem~\ref{thm:n3-two} appeals to the halting state of the algorithm of
Theorem~\ref{thm:n3-twyz}, not merely to that theorem's statement. The hypotheses of Lemma~\ref{lem:n3-stopping} do
not follow from ``envy-free partial allocation with two unassigned chores, each
of unit marginal on its holder's own bundle''. Strong connectivity of $G$ is
available only because two is the \emph{maximum} number of chores the algorithm
may leave at $n = 3$, which is what forces the component selected by (R3) to be
the whole agent set.
\end{remark}

% --------------------------------------------------------------------------
\subsection{Proof of the main theorem}
\label{sec:n3-assembly}

\begin{proof}[Proof of Theorem~\ref{thm:n3-main}]
Run the algorithm of Theorem~\ref{thm:n3-twyz} on the three cost functions
$\cost_1,\cost_2,\cost_3$, obtaining an envy-free partial allocation
$(\bundle1,\bundle2,\bundle3)$ with residue $R$ of size at most $2$.

If $\abs R = 0$ the allocation is complete and, being envy-free,
$(\alloc, \mathbf{0})$ is an envy-free solution. If $\abs R = 1$, apply
Theorem~\ref{thm:n3-one}. If $\abs R = 2$, apply Theorem~\ref{thm:n3-two}. In
each case the resulting complete allocation admits an envy-free solution with
$\optsubsidy \in \set{0,1}^3$, and Proposition~\ref{prop:zero-coordinate}
bounds the total by $2$.

For the running time, the algorithm of Theorem~\ref{thm:n3-twyz} is polynomial
and returns $(\bundle1,\bundle2,\bundle3)$ and $R$. The completion places at
most two chores into bundles and then computes a minimum-cost matching between
three agents and three bundles, 



% which the Hungarian method performs in constant
% time after the nine oracle calls determining $\cost_a(Y_j)$. 
Hence the whole
procedure is polynomial.
\end{proof}
